\section{Related Work}
\label{sec:related}

\paragraph{Catastrophic forgetting in LLMs.}
Catastrophic forgetting was first characterised in shallow connectionist networks \citep{mccloskey1989catastrophic, french1999catastrophic} and re-emerged as a central problem for continual learning in deep models \citep{kirkpatrick2017ewc, aljundi2018memory, lopez2017gem}. In large language models, sequential instruction-tuning has been shown to erode prior task performance and safety alignment \citep[e.g.,][]{luo2023empirical, shi2024sdft, wang2025llmcl}. SDFT \citep{shi2024sdft} mitigates forgetting by self-distilling the base model's outputs during fine-tuning. \citet{spurious2025} argue that some reported forgetting is \textit{spurious} --- driven by changes in the prompt-response interface rather than weight overwrite --- and motivate the counterfactual-based attribution we adopt. We differ from this line by providing a \textit{per-component}, causally-traced account of forgetting rather than a global, behavioural one.

\paragraph{Parameter-efficient fine-tuning.}
LoRA \citep{hu2022lora} established the low-rank adaptation paradigm. Subsequent work has explored adaptive rank allocation (AdaLoRA, \citealp{zhang2023adalora}), weight decomposition (DoRA, \citealp{liu2024dora}), vector-based reparameterisation (VeRA, \citealp{kopiczko2024vera}), and mixture-of-experts routing at adapter granularity \citep{mole2025}. These methods optimise \textit{how} to update a fixed target set; \tglora{} instead optimises \textit{where} to update the target set, using the forgetting topology as its selection rule. We evaluate DoRA and VeRA as baselines in \S\ref{sec:experiments} to ensure the ``matches standard LoRA'' claim is against a competitive frontier rather than the original LoRA alone.

\paragraph{Localisation-based fine-tuning.} The most direct methodological neighbour is \textbf{LOFIT} \citep{yin2024lofit}, which localises LoRA updates to a small subset of attention heads selected by \textit{new-task adaptation leverage}. \tglora{} shares the localisation mechanism but inverts the selection objective: it picks components by \textit{prior-task resistance}. LOFIT and \tglora{} thus point at overlapping but distinct target sets; we quantify the overlap via Jaccard index in \S\ref{sec:results:tglora} and treat LOFIT as a baseline. Layer-wise learning-rate scheduling \citep{lrdistribution2024} is a coarser-grained variant of the same idea (``treat components differently''), which we contrast by operating at the projection granularity.

\paragraph{Mechanistic interpretability.}
Our methodology builds on causal mediation analysis and activation patching \citep{vig2020causal, meng2022locating, meng2022memit, geva2023dissecting, wang2023ioi}. Where ROME and MEMIT localise \textit{knowledge}, we localise \textit{forgetting vulnerability}: a related but distinct object. Circuit discovery \citep{olsson2022incontext, elhage2021mathematical} and the recent open-problems agenda \citep{opmechinterp2025} provide the conceptual precedent for component-level attribution. Our contribution here is narrower than ``first per-component topology'' --- \citet{anon2026mechforget} concurrently characterises forgetting at component level via \textit{gradient} attribution. We claim the first \textit{causal-intervention} forgetting topology, validated across a second model family, and empirically compare the two attribution families in \S\ref{sec:results:ablations}.

\paragraph{Positioning.}
We unify three literatures: continual learning (what to keep), PEFT (how to update cheaply), and mechanistic interpretability (what the components do). To our knowledge, no prior work has (i) produced a causal-intervention forgetting topology validated across model families, or (ii) used such a topology to guide PEFT target selection and reported end-to-end forgetting reductions against contemporary baselines (DoRA, VeRA, LOFIT).
